% 02_related_work.tex
% Phụ trách: Thành viên 1
% ==============================================================================

\section{Related Work}
\label{sec:related-work}

\subsection{Chunking Strategies}

Chunking quyết định đơn vị văn bản được embedding và truy xuất trong hệ thống RAG \citep{lewis-etal-2020-rag}. Các phương pháp hiện có có thể được phân thành năm nhóm chính.

\paragraph{Fixed-size chunking.}
Phương pháp này chia văn bản thành các cửa sổ có số token hoặc ký tự cố định, thường kèm overlap để giảm mất thông tin tại ranh giới. Fixed-size dễ triển khai và kiểm soát chi phí, nhưng không xét đến cấu trúc hoặc tính liên kết ngữ nghĩa của nội dung \citep{qu-etal-2025-semantic}.

\paragraph{Sentence- and paragraph-based chunking.}
Nhóm phương pháp này sử dụng câu hoặc đoạn văn làm ranh giới tự nhiên. Cách tiếp cận này hạn chế việc cắt giữa câu nhưng có thể tạo ra các chunk có kích thước không đồng đều hoặc tách rời những phần nội dung phụ thuộc lẫn nhau \citep{qu-etal-2025-semantic}.

\paragraph{Semantic chunking.}
Semantic chunking xác định ranh giới dựa trên mức độ tương đồng embedding giữa các câu liên tiếp. Mục tiêu là duy trì tính liên kết ngữ nghĩa, nhưng phương pháp này làm tăng chi phí xử lý và chưa cho thấy lợi ích nhất quán so với fixed-size chunking \citep{qu-etal-2025-semantic}.

\paragraph{Structure-aware chunking.}
Structure-aware chunking khai thác heading hierarchy, paragraph, list, table, code block hoặc abstract syntax tree (AST). Các thành phần có quan hệ được tổ chức theo cấu trúc phân cấp nhằm bảo toàn ngữ cảnh của tài liệu \citep{jain-etal-2025-autochunker}.

\paragraph{LLM-assisted chunking.}
LLM-assisted hoặc prompt-based chunking sử dụng language model để xác định các đơn vị nội dung và đề xuất ranh giới linh hoạt. Các phương pháp như PIC và AutoChunker cho thấy khả năng kết hợp tín hiệu ngữ nghĩa với cấu trúc tài liệu, nhưng yêu cầu thêm chi phí suy luận và cơ chế kiểm soát đầu ra \citep{wang-etal-2025-document,jain-etal-2025-autochunker}.

\subsection{Evaluation of RAG Systems}

Đánh giá RAG thường được thực hiện ở nhiều cấp độ \citep{yu-etal-2024-rag-evaluation}. Ở cấp retrieval, Hit@\(K\), Recall@\(K\) và Mean Reciprocal Rank (MRR) đo khả năng truy xuất và xếp hạng nguồn liên quan. Ở cấp evidence, Evidence Coverage và All-Evidence-Retrieved Rate phản ánh mức độ thu hồi các bằng chứng cần thiết. Ở cấp câu trả lời, các chỉ số như Token F1, BLEU và ROUGE đo mức độ tương đồng với đáp án tham chiếu; human evaluation có thể bổ sung đánh giá về tính chính xác và đầy đủ. Việc phân biệt ba cấp độ là cần thiết vì cải thiện retrieval hoặc evidence recovery không nhất thiết dẫn đến chất lượng câu trả lời cao hơn \citep{qu-etal-2025-semantic}.

\subsection{Research Gap}

Các nghiên cứu gần đây đã so sánh fixed-size với semantic chunking hoặc đề xuất phương pháp phân đoạn mới \citep{qu-etal-2025-semantic,wang-etal-2025-document}. Tuy nhiên, tác động của việc bảo toàn cấu trúc tài liệu đến retrieval ranking, evidence recovery và answer quality vẫn chưa được đánh giá đầy đủ trong cùng một thiết lập trên tài liệu kỹ thuật có cấu trúc. Nghiên cứu này giải quyết khoảng trống đó bằng cách so sánh fixed-size, structure-aware và prompt-based chunking trong một pipeline được kiểm soát, đồng thời đánh giá riêng biệt cả ba cấp độ.